\documentclass[10pt]{article}
\usepackage[margin=12mm]{geometry}
\usepackage{graphicx}
\usepackage{caption}
\pagestyle{empty}
\captionsetup{font=small,labelfont=bf}

\begin{document}
\begin{figure}[p]
\centering
\includegraphics[width=0.92\textwidth]{fig7A_ce_full_tree_layerwise_rounds.pdf}
\vspace{-2mm}
\includegraphics[width=0.92\textwidth]{fig7B_ce_full_tree_layerwise_aggregate.pdf}
\caption{\textbf{Full-tree Pareto structure emerges across bottom-up
assignment rounds.} The measured \emph{C.\@ elegans} protein lineage was
represented as a four-root forest containing 500 internal and 504 terminal
cells, with 1,000 evaluated parent--child edges. (A) Starting from the terminal
cells, each round assigns the current child identities to the multiset of their
natural parent slots by minimum-cost bipartite matching over 301 weightings of
travel and protein-state distance. The asynchronous contraction rounds contain
504, 230, 126, 68, 38, 19, 10, and 5 edges; because measured branches terminate
at different lineage depths, rounds are bottom-up wavefronts rather than uniform
developmental layers. Each objective is expressed in standard deviations of
1,000 random permutations of the same round's parent slots and translated so
the natural assignment (black cross) is at the origin. Blue encodes retention
of the biological parent identity, so exchanges between duplicate slots of one
parent remain retained. Outlined circles mark the maximum-retention sampled
assignment. (B) Summing the eight independently optimized rounds at the same
weight produces the aggregate full-tree front. Its matched null (inset) sums
independent parent-slot permutations across rounds. The natural lineage is
sampled-dominated in the aggregate assignment space; the maximum-retention
aggregate preserves 49.9\% of natural edges.}
\label{fig:ce_full_tree_layerwise}
\end{figure}
\end{document}
